\section{Renormalization of Non-Local Operators in Dimensional Regularization}

Now let us consider the renormalization of the non-local operator, such as in Eq.~(\ref{Eq:nonlocalop}), which appears
in the LaMET approach to parton physics. In the effective theory, the non-local operator
becomes a product of local composite operators,
\begin{equation}
 O(z_2, z_1) = \overline{j}(z_2) j(z_1),
\end{equation}
after we replace the non-local Wilson line by the product of two auxiliary ``heavy quark'' fields.
 This operator can be multiplicatively renormalized by
 \begin{equation}
        O(z_2,z_1)= Z_{\bar j}Z_{j} O_R(z_2,z_1)
\end{equation}
to all orders in perturbation theory with $Z_{j}=Z_q^{1/2} Z_Q^{1/2} Z_{V_j}$, where $Z_q, Z_Q, Z_{V_j}$ are the renormalization constant for the light-quark, heavy-quark and vertex, respectively.

Thus the renormalized operator becomes $O_R=Z_{\bar j}^{-1}Z_{j}^{-1}O(z_2,z_1)$. After integration over the ``heavy quark'' field, we have
\begin{equation}\label{Eq:Orenorm}
        O_R = Z_{\bar j}^{-1}Z_{j}^{-1} \overline{\psi}(z_2) \Gamma L(z_2,z_1)\psi(z_1),
\end{equation}
where all fields on the r.h.s. are bare fields. In Ref.~\cite{Ji:2015jwa}, it has been shown that the quasi-PDF operator requires the renormalization of virtual diagrams only, which is equivalent to the result in Eq.~(\ref{Eq:Orenorm}) that we only need renormalization of the heavy-light quark currents for the operator $O$. In particular, the one-loop renormalization factor for the quasi-PDF
\beq\label{Eq:1loopZfac}
Z_O=1+\frac{\alpha_s }{\pi\epsilon}
\eeq
is consistent with the anomalous dimension of the heavy-light quark current in Eq.~(\ref{eq:hl}). The same is true at the two-loop order~\cite{Ji:2015jwa}.
